\scanpage{363}
\mspage{102}

étant le foncteur image inverse $\varphi_{n,n'}^*$, où
\[
\varphi_{n,n'}:t_{n'}\longrightarrow t_n
\]
est le composé des $t_{n'}\mathrel{\unsure{-|\to}}t_{n'-1}\mathrel{-}\ldots\mathrel{-}t_n$. (Si on n'avait pas transitivité stricte, on définirait $\varphi_{n,n'}^*$ comme le composé des foncteurs image inverse par les flèches de transition successives $t_{n'}\mathrel{-}t_{n'-1}\mathrel{-}\cdots\mathrel{-}t_{n\unsure{-|+}1}\mathrel{-}t_n$, ce ne serait $\varphi_{n,n'}^*$ qu'à isom. can. près.)

Soit \deleted{$C_\infty$}
\[
C_\infty=\underset{n\unsure{.}}{\varinjlim}\,C_n=\underset{n}{\underline{\lim}}\,C_{t_n}
\]
la catégorie $\varinjlim$. (Elle sera surtout sympa si on a pu choisir le système $(t_n)_{n\in\mathbb N}$ de telle façon que les foncteurs de transition
\[
C_{t_n}\mathrel{-}C_{t_{n\unsure{-|+}1}}
\]
soient pleinement fidèles, et peut-être même \unsure{mono|monos}, de telle façon que les $C_n$ s'identifient à des sous-catégories pleines de leur limite inductive $C_\infty$.) Je voudrais comprendre les relations précises entre les trois genres d'objets suivants :

\textcircled{a} $C$ elle-même, \deleted{\unsure{comme}} \inserted{(comme objet de $\mathrm{Cat}$, pas de $\mathrm{Cat}/T$\unsure{.|!})} et surtout les objets de $\mathrm{Cat\,hot}$ et de $\mathrm{Hot}$ qui lui correspondent.

% archive sheet number: 362
